\chapter{Confronti, cicli e membri di classe in Java}
\label{chap:confronti-cicli-java}

Le strutture condizionali introdotte nella
\cref{sec:fi-java-classi-condizioni} permettono di scegliere fra
comportamenti alternativi. Per utilizzarle correttamente è necessario
comprendere come Java confronta stringhe e oggetti, come combina
condizioni booleane e come rappresenta scelte fra più possibilità.
Le stesse condizioni governano le strutture iterative, che ripetono
un'elaborazione fino al raggiungimento di un risultato. La trattazione
si completa con i membri statici e con le regole che determinano
visibilità e durata delle variabili.
\footnote{Materiale didattico di riferimento: S. Canazza,
\emph{Java: confronti, cicli},
\nolinkurl{2026_FI_10_Java_Confronti_cicli.key.pdf}.}

\section{Confronto tra stringhe}
\label{sec:fi-java-cicli-stringhe}

In Java una stringa è un oggetto della classe \texttt{String},
non un valore di tipo primitivo (\cref{sec:fi-java-stringhe-operazioni}).
Occorre perciò distinguere il confronto dei riferimenti dal confronto
del contenuto degli oggetti.

\subsection{Uguaglianza e identità}
\label{sec:fi-java-cicli-stringhe-uguaglianza}

L'operatore \texttt{==}, applicato a due variabili di tipo riferimento,
verifica se entrambe si riferiscono allo stesso oggetto; non confronta
i caratteri delle stringhe. Il metodo \texttt{equals()}, invece,
confronta il loro contenuto. Il metodo \texttt{equalsIgnoreCase()}
confronta le stringhe senza distinguere maiuscole e minuscole, secondo
le regole della classe \texttt{String}.

\begin{unipdcode}[language=Java]
String s1 = "Stringa";
String s2 = s1;
String s3 = "String";
s3 = s3 + "a";

System.out.println(s1 == s2);       // true
System.out.println(s1 == s3);       // false
System.out.println(s1.equals(s3));  // true
System.out.println("Casa".equalsIgnoreCase("casa")); // true
\end{unipdcode}

L'assegnamento a \texttt{s2} copia il riferimento contenuto in
\texttt{s1}, mentre la concatenazione produce una nuova stringa.
Il risultato non deve essere dedotto dalla sola somiglianza dei testi:
due riferimenti possono indicare oggetti distinti contenenti gli stessi
caratteri. La \cref{fig:fi-java-cicli-stringhe-riferimenti} rappresenta
la situazione dopo le assegnazioni.

\begin{figure}[htbp]
  \centering
  \begin{tikzpicture}[unipd/diagram,x=1cm,y=1cm]
  \node[unipd/node,minimum width=1.25cm] (s1) at (-3.55,1.3) {\texttt{s1}};
  \node[unipd/node,minimum width=1.25cm] (s2) at (-3.55,0.25) {\texttt{s2}};
  \node[unipd/node,minimum width=1.25cm] (s3) at (-3.55,-1.35) {\texttt{s3}};
  \node[unipd/node,minimum width=3.0cm,minimum height=0.85cm,align=center]
    (o1) at (1.1,0.8) {\texttt{String}\\contenuto: \texttt{Stringa}};
  \node[unipd/node,minimum width=3.0cm,minimum height=0.85cm,align=center]
    (o2) at (1.1,-1.35) {\texttt{String}\\contenuto: \texttt{Stringa}};
  \draw[unipd/directed edge] (s1.east) -- ++(0.9,0) |- ($(o1.west)+(0,0.18)$);
  \draw[unipd/directed edge] (s2.east) -- ++(0.9,0) |- ($(o1.west)+(0,-0.18)$);
  \draw[unipd/directed edge] (s3.east) -- (o2.west);
\end{tikzpicture}

  \caption{Confronto fra identità dei riferimenti e uguaglianza del
  contenuto delle stringhe.}
  \label{fig:fi-java-cicli-stringhe-riferimenti}
\end{figure}

\begin{warning}[Confrontare stringhe con \texttt{==}]
Il confronto con \texttt{==} è sintatticamente valido, ma non esprime
l'uguaglianza testuale. In particolare, la condivisione di oggetti
\texttt{String} tramite il meccanismo di \emph{interning} può far
risultare veri alcuni confronti fra letterali: non bisogna affidarsi
a questo comportamento per confrontare contenuti.
\end{warning}

Un riferimento può contenere il valore \texttt{null}
(\cref{sec:fi-java-classi-null}). L'espressione
\texttt{s == null} è un confronto appropriato; invocare
un metodo su un riferimento nullo provoca invece un errore
durante l'esecuzione. Quando il primo testo è una costante,
si può scrivere \texttt{"Q".equals(s)}: il confronto
restituisce \texttt{false} anche se \texttt{s} è nullo.
Il metodo \texttt{equalsIgnoreCase()} non definisce un ordinamento
alfabetico dipendente dalla lingua e non elimina automaticamente
accenti o altri segni diacritici.

\subsection{Ordinamento lessicografico}
\label{sec:fi-java-cicli-stringhe-ordine}

Per stabilire quale stringa preceda un'altra si usa il metodo
\texttt{compareTo()}, che restituisce un valore \texttt{int}:
negativo se il ricevente precede l'argomento, nullo se i testi
coincidono e positivo se lo segue. Conta il \emph{segno} del risultato,
non la sua grandezza.

\begin{unipdcode}[language=Java]
System.out.println("carta".compareTo("castano") < 0); // true
System.out.println("casa".compareTo("casa") == 0);    // true
System.out.println("casa".compareTo("casale") < 0);  // true
\end{unipdcode}

Il confronto procede dall'inizio, cercando la prima posizione nella
quale i caratteri differiscono. Se non ne trova e una stringa è un
prefisso proprio dell'altra, la più corta precede la più lunga. Per
\texttt{String}, l'ordinamento naturale è basato sulle unità di codice
UTF-16: non coincide necessariamente con quello di un dizionario
italiano. Per esempio, le maiuscole e le minuscole sono distinte, e
anche cifre, spazi e simboli partecipano al confronto. L'ordinamento
secondo le convenzioni linguistiche richiede strumenti diversi, che
non sono necessari per i confronti introdotti qui.

\section{Confronto tra oggetti}
\label{sec:fi-java-cicli-oggetti}

La distinzione vista per le stringhe vale per tutti i riferimenti:
\texttt{a == b} verifica se i due riferimenti identificano lo stesso
oggetto (oppure sono entrambi \texttt{null}). Il metodo
\texttt{equals()} è disponibile per gli oggetti perché è dichiarato
nella classe \texttt{Object}, dalla quale derivano le classi Java.
Nella sua implementazione di base, esso verifica l'identità. Una
classe può \emph{ridefinirlo} per esprimere un'uguaglianza basata
su caratteristiche significative degli oggetti, come avviene per
\texttt{String}.

\begin{definition}[Uguaglianza per identità e per valore]
\label{def:fi-java-cicli-uguaglianza-oggetti}
L'\emph{uguaglianza per identità} richiede che due riferimenti
indichino lo stesso oggetto. L'\emph{uguaglianza per valore}, quando
prevista dalla classe, verifica invece che gli oggetti soddisfino
un criterio di equivalenza basato sul loro contenuto o stato.
\end{definition}

Non si può quindi assumere che \texttt{equals()} confronti sempre
tutti i campi: il significato del metodo dipende dalla classe e dal
contratto che essa documenta. Analogamente, \texttt{compareTo()}
non è disponibile per qualsiasi oggetto. Le classi che definiscono
un ordinamento naturale, tramite l'interfaccia \texttt{Comparable},
possono fornire questo metodo. Un risultato nullo indica equivalenza
\emph{secondo il criterio di ordinamento}, che non coincide
necessariamente con \texttt{equals()}. Per le classi create nel
\cref{chap:realizzazione-classi-java} non è stato ancora definito
un criterio generale di uguaglianza o ordinamento.

\section{Alternative multiple e istruzione \texttt{switch}}
\label{sec:fi-java-cicli-alternative}

\subsection{Catene condizionali e diramazioni annidate}
\label{sec:fi-java-cicli-if-catene}

Quando le alternative sono più di due, si può scrivere una catena
\texttt{if}/\texttt{else if}/\texttt{else}. Le condizioni vengono
esaminate nell'ordine; viene eseguito il corpo della prima condizione
vera, oppure quello finale associato a \texttt{else} se nessuna
condizione è vera.

\begin{unipdcode}[language=Java]
if (richter >= 8) {
    System.out.println("Terremoto molto forte");
} else if (richter >= 6) {
    System.out.println("Terremoto forte");
} else if (richter >= 4) {
    System.out.println("Terremoto medio");
} else if (richter >= 2) {
    System.out.println("Terremoto debole");
} else if (richter >= 0) {
    System.out.println("Terremoto molto debole");
} else {
    System.out.println("Valore negativo non valido");
}
\end{unipdcode}

La classificazione è didattica: si assume un valore finito e si
considerano non validi quelli negativi. Le soglie sono in ordine
decrescente: se \texttt{richter} vale \(7\), la prima condizione vera è
\texttt{richter >= 6}. Invertire l'ordine farebbe intercettare
ogni valore non negativo dal controllo \texttt{richter >= 0},
rendendo irraggiungibili gli altri rami. Scrivere più \texttt{if}
indipendenti, invece, potrebbe eseguire più corpi per lo stesso valore.

Le \emph{diramazioni annidate} servono quando una decisione dipende
da un'altra. Per esempio, uno scaglione di imposta può essere scelto
prima in base a una categoria e poi in base al reddito. L'istruzione
\texttt{else} viene associata sintatticamente all'\texttt{if} più
vicino che non possiede già una clausola \texttt{else}: è il problema
dell'\emph{else sospeso}. L'indentazione non modifica questa regola.
Racchiudere i rami tra graffe rende esplicita l'associazione.

\begin{unipdcode}[language=Java]
if (richter >= 0) {
    if (richter <= 4) {
        System.out.println("Terremoto leggero");
    }
} else {
    System.out.println("Valore negativo non valido");
}
\end{unipdcode}

\clearpage

\subsection{Istruzione \texttt{switch} tradizionale}
\label{sec:fi-java-cicli-switch}

L'istruzione \texttt{switch} seleziona un ramo in base al valore di
un'espressione. È particolarmente adatta a confronti con valori
discreti e distinti; non sostituisce una catena di intervalli
numerici, come quella dell'esempio precedente. Nella forma classica
si impiegano \texttt{case}, l'eventuale ramo \texttt{default} e,
normalmente, \texttt{break} per uscire dall'istruzione.

\begin{unipdcode}[language=Java]
int giorno = 3;
switch (giorno) {
    case 1:
        System.out.println("Lunedi");
        break;
    case 2:
        System.out.println("Martedi");
        break;
    case 3:
        System.out.println("Mercoledi");
        break;
    default:
        System.out.println("Altro giorno");
}
\end{unipdcode}

Le etichette \texttt{case} rappresentano valori compatibili con il
tipo dell'espressione. Lo \texttt{switch} tradizionale accetta, fra
gli altri, \texttt{byte}, \texttt{short}, \texttt{char},
\texttt{int}, le corrispondenti classi involucro, \texttt{String}
(a partire da Java 7) e i tipi enumerativi introdotti nella
\cref{sec:fi-java-cicli-enum}. Non supporta direttamente
\texttt{long} o \texttt{double} come selettori.

\begin{warning}[Caduta tra i casi]
Se un ramo raggiunge il successivo senza incontrare \texttt{break}
(o un'altra istruzione che interrompa quel percorso), l'esecuzione
\emph{prosegue} nel caso seguente, senza verificare nuovamente le
etichette. Questo comportamento, detto \emph{fall-through}, è talvolta
voluto per raggruppare più valori, ma può introdurre errori.
\end{warning}

\begin{unipdcode}[language=Java]
switch (giorno) {
    case 1:
    case 2:
    case 3:
    case 4:
    case 5:
        System.out.println("Giorno lavorativo");
        break;
    case 6:
    case 7:
        System.out.println("Fine settimana");
        break;
    default:
        System.out.println("Giorno non valido");
}
\end{unipdcode}

L'assenza dei primi quattro \texttt{break} è intenzionale: tutte
le relative etichette conducono allo stesso corpo. Al contrario,
omettere il \texttt{break} dopo la stampa di ``Giorno lavorativo''
provocherebbe anche la stampa di ``Fine settimana''.

\clearpage

Lo stesso selettore può essere una stringa. Per esempio, la
scrittura seguente distingue due risposte testuali; qui le
etichette \texttt{case} effettuano un confronto per contenuto,
non per identità dei riferimenti:

\begin{unipdcode}[language=Java]
String risposta = "si";
switch (risposta) {
    case "si":
        System.out.println("Confermato");
        break;
    case "no":
        System.out.println("Rifiutato");
        break;
    default:
        System.out.println("Risposta non riconosciuta");
}
\end{unipdcode}

\subsection{Espressioni \texttt{switch} e enumerazioni}
\label{sec:fi-java-cicli-switch-expression}

A partire da Java 14, lo \texttt{switch} può essere utilizzato
anche come \emph{espressione} che produce un valore. Le etichette
con la freccia \texttt{->} delimitano un singolo ramo, senza
\emph{fall-through}; il compilatore richiede che ogni possibile
valore del selettore sia gestito. Nell'esempio, il ramo
\texttt{default} copre gli interi diversi da quelli elencati.

\begin{unipdcode}[language=Java]
int giorno = 3;
String tipoGiorno = switch (giorno) {
    case 1, 2, 3, 4, 5 -> "Giorno lavorativo";
    case 6, 7 -> "Fine settimana";
    default -> "Giorno non valido";
};
\end{unipdcode}

È importante distinguere la versione \emph{istruzione}, utilizzata
per eseguire azioni, dalla versione \emph{espressione}, impiegata
per calcolare un risultato. Le etichette con freccia possono essere
usate anche in istruzioni \texttt{switch} nelle versioni moderne
di Java. Quando il selettore è un \texttt{enum}, i casi possono
indicare direttamente le costanti senza ripetere il nome del tipo;
la definizione e l'esempio completo sono nella
\cref{sec:fi-java-cicli-enum}.

\section{Il tipo \texttt{boolean} e le espressioni logiche}
\label{sec:fi-java-cicli-boolean}

\begin{definition}[Tipo booleano]
\label{def:fi-java-cicli-boolean}
Il tipo primitivo \texttt{boolean} ha esattamente due valori,
\texttt{true} e \texttt{false}. Le condizioni degli enunciati
\texttt{if} e dei cicli devono essere espressioni di tipo
\texttt{boolean}; gli interi non sono convertiti implicitamente
in valori logici.
\end{definition}

Per esempio, \texttt{x > 0} è un'espressione relazionale il cui
risultato può essere assegnato a una variabile:

\begin{unipdcode}[language=Java]
boolean positivo = x > 0;
if (positivo) {
    System.out.println("Valore positivo");
}
\end{unipdcode}

Quando una variabile booleana memorizza lo stato di una condizione
utile al controllo del programma, viene spesso chiamata
\emph{flag}. I metodi che restituiscono valori booleani sono detti
\emph{predicativi}: \texttt{Character.isDigit(ch)} verifica, per
esempio, se un carattere è una cifra, mentre
\texttt{scanner.hasNextInt()} verifica se il prossimo elemento
dell'ingresso può essere letto come \texttt{int}
(\cref{sec:fi-java-scanner-numeri}). I nomi dei predicati spesso
iniziano con \texttt{is} oppure \texttt{has}.

\subsection{Operatori logici e cortocircuito}
\label{sec:fi-java-cicli-logici}

Gli operatori \texttt{\&\&}, \texttt{||} e \texttt{!}
rappresentano rispettivamente congiunzione, disgiunzione e
negazione. Le corrispondenti tabelle di verità sono riassunte
nella \cref{tab:fi-java-cicli-booleani}.

\begin{table}[htbp]
\centering
\begin{tabular}{ccccc}
\toprule
\(A\) & \(B\) & \texttt{A \&\& B} & \texttt{A || B} & \texttt{!A} \\
\midrule
\texttt{true} & \texttt{true} & \texttt{true} & \texttt{true} & \texttt{false} \\
\texttt{true} & \texttt{false} & \texttt{false} & \texttt{true} & \texttt{false} \\
\texttt{false} & \texttt{true} & \texttt{false} & \texttt{true} & \texttt{true} \\
\texttt{false} & \texttt{false} & \texttt{false} & \texttt{false} & \texttt{true} \\
\bottomrule
\end{tabular}
\caption{Valori degli operatori logici su due condizioni booleane.}
\label{tab:fi-java-cicli-booleani}
\end{table}

Con \texttt{\&\&} e \texttt{||}, Java utilizza la
\emph{valutazione a cortocircuito}: il secondo operando non viene
valutato quando il primo determina già il risultato. Nel seguente
controllo, la chiamata a \texttt{length()} avviene soltanto se
\texttt{s} non è nullo:

\begin{unipdcode}[language=Java]
if (s != null && s.length() > 0) {
    System.out.println("Stringa non vuota");
}
\end{unipdcode}

Invertendo le condizioni, la protezione dal riferimento nullo
verrebbe meno e potrebbe verificarsi una
\texttt{NullPointerException}. Il cortocircuito riguarda quindi
anche la correttezza, non soltanto l'efficienza.

\subsection{Precedenze e confronti errati}
\label{sec:fi-java-cicli-precedenze}

Gli operatori aritmetici e quelli relazionali precedono in generale
le combinazioni logiche; tra i tre operatori logici considerati,
\texttt{!} ha precedenza maggiore di \texttt{\&\&}, che precede
\texttt{||}. Le parentesi chiariscono o modificano il
raggruppamento, mentre le regole di associatività risolvono i casi
di operatori con la stessa precedenza. Per esempio,
\texttt{2 * x \% 2} viene raggruppato come
\texttt{(2 * x) \% 2}.

\begin{unipdcode}[language=Java]
if (!(x < 0 || x > 10)) {
    System.out.println("Nell'intervallo [0, 10]");
}
\end{unipdcode}

Il tentativo di scrivere \texttt{!x < 0} per indicare \(x\ge0\)
non è corretto: \texttt{!} richiede un operando booleano e non
può essere applicato a un numero. Analogamente, non si possono
concatenare confronti come nelle disuguaglianze matematiche:

\noindent
\begin{minipage}[t]{0.48\linewidth}
\vspace{0pt}
\textbf{Confronto di intervallo}

\begin{unipdcode}[language=Java,
  basicstyle=\ttfamily\scriptsize,
  numbers=none,
  xleftmargin=1em,
  xrightmargin=1em,
  framexleftmargin=0.3em]
// NON valido: 0 <= x <= 1
if (0 <= x && x <= 1) {
    System.out.println("Tra zero e uno");
}
\end{unipdcode}
\end{minipage}%
\hfill
\begin{minipage}[t]{0.48\linewidth}
\vspace{0pt}
\textbf{Confronto di valori}

\begin{unipdcode}[language=Java,
  basicstyle=\ttfamily\scriptsize,
  numbers=none,
  xleftmargin=1em,
  xrightmargin=1em,
  framexleftmargin=0.3em]
// NON valido: x && y > 0
if (x > 0 && y > 0) {
    System.out.println("Entrambi positivi");
}
\end{unipdcode}
\end{minipage}

L'espressione \texttt{0 <= x} produce un valore booleano, che
non può essere confrontato numericamente con \texttt{1}; inoltre,
\texttt{\&\&} richiede due operandi booleani. Questi errori
vengono individuati dal compilatore. Il confronto fra due
numeri in virgola mobile presenta un problema differente,
relativo alla precisione, già trattato nella
\cref{sec:fi-java-classi-floating}.

\section{Tipi enumerativi \texttt{enum}}
\label{sec:fi-java-cicli-enum}

Quando una variabile deve rappresentare uno fra un numero finito e
prestabilito di stati, si può utilizzare un tipo \emph{enumerativo}.
Una dichiarazione \texttt{enum} introduce in Java un tipo speciale
di classe, le cui istanze sono le costanti elencate nella
dichiarazione. Per esempio:

\begin{unipdcode}[language=Java]
enum Stagione {
    PRIMAVERA, ESTATE, AUTUNNO, INVERNO
}
\end{unipdcode}

All'interno di un metodo si può dichiarare una variabile di questo
tipo e assegnarle una delle costanti:
\begin{unipdcode}[language=Java]
Stagione stagione = Stagione.ESTATE;
\end{unipdcode}

La variabile può contenere uno dei valori dichiarati oppure
\texttt{null}, trattandosi di un tipo riferimento. Non è possibile
assegnarle arbitrariamente una stringa o una costante di un altro
\texttt{enum}. Le costanti di una stessa enumerazione sono oggetti
unici, perciò è corretto confrontarle mediante \texttt{==}.

\begin{example}[Selezione con un tipo enumerativo]
Il seguente esempio seleziona una costante di tipo
\texttt{Stagione} mediante \texttt{switch}. Le etichette
\texttt{case} utilizzano i nomi delle costanti senza il prefisso
\texttt{Stagione.}.
\begin{unipdcode}[language=Java]
enum Stagione {
    PRIMAVERA, ESTATE, AUTUNNO, INVERNO
}

public class StagioniTester {
    public static void main(String[] args) {
        Stagione stagione = Stagione.ESTATE;
        switch (stagione) {
            case PRIMAVERA:
                System.out.println("Primavera");
                break;
            case ESTATE:
                System.out.println("Estate");
                break;
            case AUTUNNO:
                System.out.println("Autunno");
                break;
            case INVERNO:
                System.out.println("Inverno");
                break;
        }
    }
}
\end{unipdcode}

\end{example}

Un tipo enumerativo può anche dichiarare campi, costruttori e
metodi. In questo caso ogni costante viene costruita con i valori
previsti dalla dichiarazione; non è possibile crearne liberamente
altre con \texttt{new}. Le classi enumerative non possono essere
estese mediante ereditarietà. Queste proprietà rendono
\texttt{enum} adatto a modellare un insieme chiuso di valori,
non un catalogo che deve cambiare dinamicamente durante
l'esecuzione.

\begin{example}[Enumerazioni con dati e metodi]
Nella seguente enumerazione ogni valore possiede una descrizione,
assegnata al momento della costruzione delle costanti e restituita
da un metodo di istanza. Si tratta di un esempio ridotto del
meccanismo utilizzato per associare proprietà a categorie fisse.
\begin{unipdcode}[language=Java]
enum Priorita {
    BASSA("Da pianificare"),
    ALTA("Da affrontare subito");

    private final String descrizione;

    Priorita(String descrizione) {
        this.descrizione = descrizione;
    }

    public String getDescrizione() {
        return descrizione;
    }
}

public class PrioritaTester {
    public static void main(String[] args) {
        System.out.println(Priorita.ALTA.getDescrizione());
    }
}
\end{unipdcode}

\end{example}

I principali vantaggi sono la \emph{sicurezza del tipo}, la
leggibilità dei nomi e l'integrazione con \texttt{switch};
evitano, per esempio, di confondere codici numerici attribuiti
a categorie diverse. Il limite è che l'insieme delle costanti
è fissato nel codice e una sua modifica richiede di aggiornare
il programma. Inoltre, convertire dati testuali esterni in
costanti richiede un'operazione esplicita e la gestione di
valori non riconosciuti. Le enumerazioni, dunque, non
sostituiscono le strutture dati per insiemi di valori variabili:
la distinzione tra specifica di un tipo e rappresentazione dei
dati richiama il concetto di tipo astratto introdotto nella
\cref{sec:fi-adt-strutture}.

\section{Il ciclo \texttt{while}}
\label{sec:fi-java-cicli-while}

Le strutture iterative eseguono più volte un corpo di istruzioni.
Nel ciclo \texttt{while} la condizione viene verificata
\emph{prima} di ogni iterazione. Se è falsa fin dall'inizio,
il corpo non viene eseguito neppure una volta. La sua forma
schematica è:

\begin{unipdcode}[language=Java]
while (condizione) {
    // corpo del ciclo
}
\end{unipdcode}

La condizione deve avere tipo \texttt{boolean}, e il corpo può
essere un singolo enunciato o un blocco. La
\cref{fig:fi-java-cicli-controllo} mette a confronto questo
schema con quello del ciclo \texttt{do}/\texttt{while},
introdotto più avanti.

\begin{figure}[htbp]
  \centering
  \begin{tikzpicture}[unipd/diagram,x=1cm,y=1cm]
  \node at (-3.1,2.35) {\textbf{\texttt{while}}};
  \node at (3.0,2.35) {\textbf{\texttt{do}/\texttt{while}}};
  \node[unipd/node,minimum width=2.05cm,minimum height=0.65cm] (wtest) at (-3.1,1.25) {Condizione?};
  \node[unipd/node,minimum width=2.05cm,minimum height=0.65cm] (wbody) at (-3.1,-0.1) {Corpo};
  \node[unipd/node,minimum width=1.30cm] (wend) at (-3.1,-1.75) {Fine};
  \node[unipd/node,minimum width=2.05cm,minimum height=0.65cm] (dbody) at (3.0,1.25) {Corpo};
  \node[unipd/node,minimum width=2.05cm,minimum height=0.65cm] (dtest) at (3.0,-0.1) {Condizione?};
  \node[unipd/node,minimum width=1.30cm] (dend) at (3.0,-1.75) {Fine};
  \draw[unipd/directed edge] (wtest.south) -- node[right]{vero} (wbody.north);
  \draw[unipd/directed edge] (wbody.west) -- ++(-1.05,0) |- (wtest.west);
  \draw[unipd/directed edge] (wtest.east) -- ++(1.05,0) |- node[pos=0.15,right]{falso} (wend.east);
  \draw[unipd/directed edge] (dbody.south) -- (dtest.north);
  \draw[unipd/directed edge] (dtest.east) -- ++(1.05,0) |- node[pos=0.10,right]{vero} (dbody.east);
  \draw[unipd/directed edge] (dtest.south) -- node[right]{falso} (dend.north);
\end{tikzpicture}

  \caption{Verifica preliminare nel ciclo \texttt{while} e
  verifica successiva nel ciclo \texttt{do}/\texttt{while}.}
  \label{fig:fi-java-cicli-controllo}
\end{figure}

\subsection{Capitalizzazione annuale}
\label{sec:fi-java-cicli-investimento}

Si consideri un conto con saldo iniziale di \(20\,000\) euro e
tasso annuo costante del \(5\%\), con capitalizzazione degli
interessi alla fine di ogni anno. Si vuole trovare il minimo
numero di anni necessario affinché il saldo sia almeno il doppio
di quello iniziale. Indicando con \(B_n\) il saldo dopo \(n\)
anni, si ha
\[
  B_0=20\,000,\qquad B_{n+1}=1{,}05\,B_n.
\]
L'algoritmo mantiene il saldo corrente e il numero di anni
trascorsi; a ogni iterazione aumenta prima il contatore e poi
accredita gli interessi. La sua realizzazione può utilizzare
\texttt{BankAccount}, già definita nella
\cref{sec:fi-java-classi-incapsulamento}:

\begin{unipdcode}[language=Java]
public class SimpleInvestmentTester {
    public static void main(String[] args) {
        final double INITIAL_BALANCE = 20000;
        final double RATE = 5;
        BankAccount account = new BankAccount(INITIAL_BALANCE);
        int years = 0;

        while (account.getBalance() < 2 * INITIAL_BALANCE) {
            years++;
            double interest = account.getBalance() * RATE / 100;
            account.deposit(interest);
        }
        System.out.println("Raddoppio dopo " + years + " anni");
    }
}
\end{unipdcode}


All'inizio di ciascuna verifica della condizione, il contatore
\texttt{years} è il numero di capitalizzazioni già effettuate e
\texttt{account.getBalance()} è il saldo corrispondente. Questo
è un \emph{invariante di ciclo}, utile per giustificare il
risultato. Poiché il tasso è positivo e il saldo iniziale è
positivo, il saldo aumenta di anno in anno e supera infine
la soglia: il ciclo termina. Il minimo \(n\) cercato è
\(15\), poiché \(1{,}05^{14}<2\) e \(1{,}05^{15}>2\).

La conclusione dipende dalle ipotesi: con tasso nullo e
obiettivo superiore al saldo corrente, il ciclo non terminerebbe.
Un ciclo infinito può inoltre derivare da una condizione errata,
da un contatore mai aggiornato o da un aggiornamento nella direzione
sbagliata. La terminazione va dedotta dal comportamento dell'algoritmo,
non dal fatto che il programma supera alcuni casi di prova.

\subsection{Cicli annidati e condizioni di uscita}
\label{sec:fi-java-cicli-while-annidati}

Il corpo di un ciclo può contenere un altro ciclo. In generale,
per ogni iterazione del ciclo esterno, il ciclo interno riparte
dalle proprie condizioni iniziali e può eseguire più iterazioni.
Il numero complessivo di operazioni dipende da entrambi i cicli;
per la relativa analisi si applica il modello di costo della
\cref{chap:complessita}.

Per esempio, un ciclo esterno può selezionare tre righe e un ciclo
interno stampare quattro valori per ciascuna riga:

\begin{unipdcode}[language=Java]
int riga = 1;
while (riga <= 3) {
    int colonna = 1;
    while (colonna <= 4) {
        System.out.print(riga * colonna + " ");
        colonna++;
    }
    System.out.println();
    riga++;
}
\end{unipdcode}

Il contatore interno viene inizializzato nuovamente a ogni
iterazione esterna. Se non venisse reimpostato, dopo la prima
riga il ciclo interno non eseguirebbe ulteriori iterazioni.

Una situazione ricorrente è la lettura di una quantità non nota di
dati, terminata da un \emph{valore sentinella}. La sentinella
segnala la fine dell'ingresso e non appartiene ai dati da elaborare.
Il seguente programma somma i numeri inseriti fino allo zero, che
non viene conteggiato.

\begin{unipdcode}[language=Java]
import java.util.Scanner;

public class SommaSentinella {
    public static void main(String[] args) {
        Scanner in = new Scanner(System.in);
        int somma = 0;
        int conteggio = 0;
        boolean continua = true;

        while (continua) {
            System.out.print("Numero (0 per terminare): ");
            int numero = in.nextInt();
            if (numero == 0) {
                continua = false;
            } else {
                somma += numero;
                conteggio++;
            }
        }
        System.out.println("Numeri letti: " + conteggio);
        System.out.println("Somma: " + somma);
    }
}
\end{unipdcode}


La variabile booleana \texttt{continua} è un flag: la condizione
viene modificata dopo la lettura, impedendo di dover duplicare
l'operazione prima e dentro il ciclo. Lo stesso schema funziona
con una sentinella testuale come \texttt{"Q"}; il confronto deve
allora utilizzare \texttt{equals()} oppure
\texttt{equalsIgnoreCase()} (\cref{sec:fi-java-cicli-stringhe-uguaglianza}).
Questa struttura viene talvolta chiamata \emph{ciclo e mezzo},
perché la decisione di proseguire avviene fra acquisizione ed
elaborazione del dato.

In alternativa si può utilizzare un ciclo \texttt{while (true)}
interrotto da \texttt{break}. L'istruzione \texttt{break}
termina il ciclo più interno che la contiene; \texttt{continue}
interrompe invece l'iterazione corrente e passa alla successiva
verifica (oppure alla fase di aggiornamento nel \texttt{for}).
Non sono istruzioni intrinsecamente errate, ma il loro impiego
ripetuto può rendere meno evidente il flusso di controllo.

\subsection{Il ciclo \texttt{do}/\texttt{while}}
\label{sec:fi-java-cicli-do-while}

Quando occorre eseguire almeno una volta il corpo e verificare
\emph{dopo} se ripeterlo, si utilizza \texttt{do}/\texttt{while}.
La sintassi termina con un punto e virgola dopo la condizione:

\begin{unipdcode}[language=Java]
double rate;
do {
    System.out.println("Inserire un tasso positivo:");
    rate = in.nextDouble();
} while (rate <= 0);
\end{unipdcode}

L'esempio presuppone che \texttt{in} sia uno \texttt{Scanner}
già inizializzato e che l'ingresso contenga numeri finiti:
non gestisce valori non numerici né valori speciali come
\texttt{NaN}. Il corpo è eseguito almeno una volta e
un valore finito non positivo provoca una nuova lettura.
Rispetto a \texttt{while}, non occorre attribuire a \texttt{rate}
un valore artificiale prima di leggere il primo dato.

\section{Il ciclo \texttt{for}}
\label{sec:fi-java-cicli-for}

Quando un ciclo utilizza una variabile di controllo con
inizializzazione, condizione e aggiornamento riconoscibili,
la forma \texttt{for} riunisce questi elementi nell'intestazione:

\begin{unipdcode}[language=Java]
for (int i = 0; i < 10; i++) {
    System.out.println(i);
}
\end{unipdcode}

L'inizializzazione viene eseguita una volta sola; prima di ogni
iterazione si verifica la condizione; dopo il corpo si esegue
l'aggiornamento, poi si ripete il controllo. Un ciclo
\texttt{for} può essere espresso mediante un \texttt{while}
con lo stesso ordine delle operazioni:

\begin{unipdcode}[language=Java]
int i = 0;
while (i < 10) {
    System.out.println(i);
    i++;
}
\end{unipdcode}

L'equivalenza è immediata in assenza di istruzioni che alterino
il flusso, come \texttt{continue}: in un \texttt{for},
\texttt{continue} passa prima alla fase di aggiornamento.
L'incremento non deve necessariamente essere di un'unità;
può essere negativo e la variabile di controllo può avere anche
tipo \texttt{double}, benché in quest'ultimo caso gli errori
di arrotondamento richiedano particolare attenzione
(\cref{sec:fi-java-classi-floating}).

Una variabile dichiarata nell'inizializzazione del \texttt{for}
è visibile nell'intestazione e nel corpo, ma non dopo il ciclo.
Se il suo valore finale serve successivamente, va dichiarata
prima dell'enunciato \texttt{for}. Le regole generali di
visibilità sono precisate nella \cref{sec:fi-java-cicli-scope}.

\subsection{Cicli \texttt{for} annidati}
\label{sec:fi-java-cicli-for-annidati}

Per stampare una tabella delle potenze \(x^y\), con
\(1\le x\le4\) e \(1\le y\le5\), si possono utilizzare
due cicli: quello esterno seleziona la riga e quello interno
calcola i valori delle colonne. Per valori piccoli, interi e
non negativi, si può usare \texttt{Math.pow()} con conversione
controllata all'intero, oppure un prodotto ripetuto.

\begin{unipdcode}[language=Java]
public class TabellaPotenze {
    public static void main(String[] args) {
        for (int x = 1; x <= 4; x++) {
            for (int y = 1; y <= 5; y++) {
                int p = (int) Math.round(Math.pow(x, y));
                System.out.printf("%6d", p);
            }
            System.out.println();
        }
    }
}
\end{unipdcode}


Il programma produce quattro righe e cinque colonne. La stampa
con \texttt{printf()} mantiene le colonne allineate, come già
illustrato nella \cref{sec:fi-java-printf}. Un ulteriore ciclo
interno potrebbe costruire gli spazi di riempimento, ma non
sarebbe necessario se si utilizza la formattazione disponibile.

\subsection{Errori di controllo e leggibilità}
\label{sec:fi-java-cicli-errori}

Gli \emph{errori per scarto di uno} derivano frequentemente da
una condizione finale scelta in modo sbagliato. Se gli indici
validi di una stringa sono \(0,\ldots,n-1\), il ciclo naturale è
\texttt{for (int i = 0; i < s.length(); i++)}; utilizzare
\texttt{i <= s.length()} provocherebbe un accesso oltre la fine.
Analogamente, nel problema dell'investimento si verifica
\texttt{balance < target} perché il raggiungimento esatto della
soglia deve terminare il ciclo.

Il seguente ciclo è valido e ha un corpo vuoto:
\begin{unipdcode}[language=Java]
for (int i = 0; i < 10; i++);
\end{unipdcode}
\begin{warning}[Punto e virgola dopo \texttt{for}]
Il punto e virgola costituisce da solo il corpo del ciclo.
Le istruzioni scritte dopo non appartengono al ciclo: questa
differenza sintattica può produrre un errore logico non rilevato
dal compilatore.
\end{warning}

Condizioni di uscita come \texttt{i != n} sono fragili se
l'aggiornamento può oltrepassare \texttt{n} o muoversi nella
direzione opposta; quando possibile è preferibile usare
relazioni coerenti con la direzione del ciclo, come
\texttt{i < n}. Un altro errore frequente consiste nel modificare
la variabile di controllo sia nell'intestazione sia nel corpo,
rendendo difficile prevedere le iterazioni effettive.
Le prove sui casi limite aiutano a trovare tali errori, ma
non sostituiscono il ragionamento sulla terminazione.

\section{Metodi e variabili statiche}
\label{sec:fi-java-cicli-static}

In una classe Java i membri non statici appartengono alle
singole istanze, mentre quelli dichiarati con \texttt{static}
appartengono alla classe. La differenza non riguarda soltanto
la sintassi di invocazione, ma anche la presenza o meno di uno
stato specifico dell'oggetto ricevente.

\subsection{Metodi statici e parametro implicito}
\label{sec:fi-java-cicli-metodi-statici}

Un metodo di istanza riceve implicitamente l'oggetto sul quale
opera; al suo interno tale riferimento è indicato da
\texttt{this}, come descritto nella
\cref{sec:fi-java-classi-parametri}. Un metodo statico, invece,
non possiede un parametro implicito e può essere invocato
attraverso il nome della classe.

\begin{unipdcode}[language=Java]
public class Financial {
    public static double percentOf(double percent, double amount) {
        return percent * amount / 100;
    }
}
\end{unipdcode}

La chiamata \texttt{Financial.percentOf(5, 200)} restituisce
\texttt{10.0}, senza richiedere un oggetto \texttt{Financial}.
Le classi che raccolgono metodi di questo tipo sono spesso dette
\emph{classi di utilità}; \texttt{Math} è un esempio della
libreria standard (\cref{sec:fi-java-math}). Un metodo statico
non può utilizzare direttamente \texttt{this} né leggere un
campo di istanza senza disporre di un riferimento a un oggetto.

\subsection{Variabili di classe e costanti}
\label{sec:fi-java-cicli-campi-statici}

Un campo statico ha un'unica copia associata alla classe,
indipendentemente dal numero di oggetti costruiti. Se si vuole
assegnare un numero progressivo a ciascun nuovo conto,
il contatore dell'ultimo numero assegnato deve essere condiviso.
La \cref{fig:fi-java-cicli-statici-istanze} distingue la copia
unica del contatore dai numeri assegnati alle singole istanze.

\begin{figure}[htbp]
  \centering
  \begin{tikzpicture}[unipd/diagram,x=1cm,y=1cm]
  \node[unipd/node,minimum width=5.2cm,minimum height=0.85cm,align=center]
    (shared) at (0,1.65) {\texttt{ContoNumerato} (classe)\\\texttt{static lastAssignedNumber = 2}};
  \node[unipd/node,minimum width=3.5cm,minimum height=0.9cm,align=center]
    (a) at (-2.65,-0.55) {Oggetto \texttt{a}\\\texttt{accountNumber = 1}};
  \node[unipd/node,minimum width=3.5cm,minimum height=0.9cm,align=center]
    (b) at (2.65,-0.55) {Oggetto \texttt{b}\\\texttt{accountNumber = 2}};
  \draw[unipd/directed edge] (a.north) -- ++(0,0.45) -| ($(shared.south)+(-1.3,0)$);
  \draw[unipd/directed edge] (b.north) -- ++(0,0.45) -| ($(shared.south)+(1.3,0)$);
\end{tikzpicture}

  \caption{Un campo statico condiviso e due campi di istanza,
  ciascuno appartenente a un oggetto distinto.}
  \label{fig:fi-java-cicli-statici-istanze}
\end{figure}

\clearpage

\begin{unipdcode}[language=Java]
public class ContoNumerato {
    private static int lastAssignedNumber = 0;
    private final int accountNumber;

    public ContoNumerato() {
        lastAssignedNumber++;
        accountNumber = lastAssignedNumber;
    }

    public int getAccountNumber() {
        return accountNumber;
    }

    public static void main(String[] args) {
        ContoNumerato a = new ContoNumerato();
        ContoNumerato b = new ContoNumerato();
        System.out.println(a.getAccountNumber()); // 1
        System.out.println(b.getAccountNumber()); // 2
    }
}
\end{unipdcode}


Dopo due costruzioni consecutive, gli oggetti ricevono
rispettivamente i numeri \(1\) e \(2\). Se
\texttt{lastAssignedNumber} fosse un campo di istanza
inizializzato a zero, ogni nuovo conto ripartirebbe da \(1\).
Il campo statico viene inizializzato una sola volta durante
l'inizializzazione della classe; azzerarlo nel costruttore
annullerebbe il conteggio progressivo. Il suo valore appartiene
allo stato condiviso della classe, non a un singolo conto.

L'uso di campi statici \emph{modificabili} deve essere motivato:
le loro variazioni possono influenzare tutte le istanze e
rendere meno prevedibile il comportamento dei metodi. Sono
invece comuni le costanti \texttt{public static final}, come
\texttt{Math.PI}, che si accedono qualificandone il nome.
\texttt{final} impedisce di riassegnare il campo, mentre
\texttt{static} ne determina l'appartenenza alla classe;
nessuno dei due modificatori implica automaticamente l'altro.

\section{Ambito di visibilità, ciclo di vita e memoria}
\label{sec:fi-java-cicli-scope}

L'\emph{ambito di visibilità} (\emph{scope}) di una variabile è la
parte del codice sorgente in cui il suo nome può essere utilizzato.
Il \emph{ciclo di vita} è invece il periodo in cui la variabile
esiste durante l'esecuzione. I due concetti non coincidono: il primo
riguarda le regole di visibilità dei nomi, il secondo l'esecuzione.
Per una variabile di tipo riferimento, inoltre, il ciclo di vita
non coincide necessariamente con quello dell'oggetto referenziato.

\clearpage

\begin{table}[htbp]
\centering
\small
\begin{tabular}{p{2.45cm}p{5.2cm}p{5.0cm}}
\toprule
Categoria & Ambito di visibilità & Ciclo di vita, in generale \\
\midrule
Locale & Blocco nel quale è dichiarata, secondo le regole di Java & Durante l'esecuzione del blocco nell'invocazione corrente \\
Parametro & Corpo del metodo o costruttore ricevente & Durante l'invocazione corrispondente \\
Di istanza & Membri della classe, secondo i modificatori di accesso & Legato all'esistenza dell'oggetto \\
Statica & Membri della classe, secondo i modificatori di accesso & Legato all'inizializzazione e alla permanenza della classe nella JVM \\
\bottomrule
\end{tabular}
\caption{Le quattro categorie principali di variabili considerate nel corso.}
\label{tab:fi-java-cicli-variabili}
\end{table}

Una variabile locale dichiarata dentro un blocco \texttt{if}
non è visibile dopo la sua chiusura; una variabile dichiarata
prima del blocco continua invece a essere accessibile al suo
termine. Il compilatore Java verifica inoltre che una
variabile locale sia \emph{definitivamente assegnata} prima
di essere letta: una dichiarazione come \texttt{int n;}
è valida, ma utilizzare \texttt{n} senza averle assegnato
un valore causa un errore di compilazione. I campi di istanza
e quelli statici, al contrario, ricevono automaticamente
valori iniziali predefiniti
(\cref{sec:fi-java-classi-costruttore-predefinito}).

\subsection{Oggetti, stack, heap e raccolta della memoria}
\label{sec:fi-java-cicli-memoria}

Nel modello concettuale di esecuzione della JVM, ogni invocazione
di metodo dispone di un proprio \emph{frame} che contiene le
variabili locali e i parametri; gli oggetti sono rappresentati
nell'\emph{heap}, cioè nell'area di memoria gestita dalla JVM.
Una variabile locale di tipo riferimento contiene il riferimento
a un oggetto, non l'oggetto stesso. La
\cref{fig:fi-java-cicli-memoria} illustra questa distinzione.

\begin{figure}[htbp]
  \centering
  \begin{tikzpicture}[unipd/diagram,x=1cm,y=1cm]
  \node at (-2.8,1.25) {\textbf{Frame di metodo}};
  \node at (2.7,1.25) {\textbf{Heap (modello logico)}};
  \node[unipd/node,minimum width=2.5cm,minimum height=1.0cm,align=center]
    (r) at (-2.8,0) {Variabile locale\\\texttt{b}: riferimento};
  \node[unipd/node,minimum width=3.3cm,minimum height=1.0cm,align=center]
    (o) at (2.7,0) {Oggetto \texttt{BankAccount}\\\texttt{balance = 100.0}};
  \draw[unipd/directed edge] (r.east) -- (o.west);
\end{tikzpicture}

  \caption{Modello concettuale di un riferimento locale a un oggetto
  contenente campi di istanza.}
  \label{fig:fi-java-cicli-memoria}
\end{figure}

Questa rappresentazione è utile per comprendere il programma, ma
non descrive necessariamente la disposizione fisica della memoria:
la JVM può applicare ottimizzazioni che modificano l'allocazione
effettiva. Quando un oggetto non è più \emph{raggiungibile} da
riferimenti utilizzabili dal programma, può essere recuperato
automaticamente mediante il \emph{garbage collector} (GC).
La raccolta non avviene necessariamente subito dopo la perdita
dell'ultimo riferimento, e neppure la chiamata
\texttt{System.gc()} garantisce una raccolta immediata.

La scelta del meccanismo di raccolta influenza due caratteristiche
dell'esecuzione: il \emph{throughput}, cioè la quantità di lavoro
completata nell'unità di tempo, e la \emph{latenza}, cioè il tempo
richiesto da una singola operazione o pausa. Fra i raccoglitori
disponibili in diverse versioni e distribuzioni della JVM vi sono
\emph{Serial GC} (che esegue la raccolta mediante un solo thread), \emph{Parallel GC} (orientato al throughput),
\emph{G1} (compromesso fra lavoro e pause), \emph{ZGC} e
\emph{Shenandoah} (progettati per ridurre le pause). Nelle
versioni comuni di HotSpot, G1 è il raccoglitore predefinito a
partire da Java 9 e può essere richiesto con
\texttt{-XX:+UseG1GC}; il precedente CMS è stato rimosso
in Java 14. Disponibilità e comportamento dipendono comunque
dalla versione e dalla distribuzione impiegate. Nessun
raccoglitore elimina la necessità di evitare riferimenti
inutilmente mantenuti a oggetti non più necessari.

\clearpage

\section{Visibilità dei membri e sovrapposizione dei nomi}
\label{sec:fi-java-cicli-shadowing}

I membri di una classe sono soggetti ai modificatori di accesso
introdotti nella \cref{sec:fi-java-classi-incapsulamento}.
Un membro \texttt{private} è accessibile nel codice della
classe che lo dichiara, anche quando il metodo opera su
un'altra istanza della stessa classe. Un membro \texttt{public}
può essere utilizzato dal codice esterno nei contesti consentiti
dal linguaggio, qualificandolo con un riferimento a un oggetto,
se è un membro di istanza, oppure con il nome della classe,
se è statico.

All'interno di un metodo di istanza, i nomi dei campi e dei metodi
di istanza possono essere utilizzati senza specificare
\texttt{this}; quelli statici possono essere qualificati con il
nome della classe. Per accedere a un campo nascosto da un nome locale, per esempio
da un parametro, occorre invece qualificare il riferimento.

\begin{definition}[Shadowing]
\label{def:fi-java-cicli-shadowing}
Nel caso dei campi di una classe, si ha \emph{shadowing}
(oscuramento di un nome) quando una dichiarazione in un ambito
più interno nasconde un campo omonimo. Il nome semplice si riferisce allora
alla dichiarazione più vicina; il campo di istanza rimane
accessibile mediante \texttt{this.nome}, e quello statico
mediante \texttt{NomeClasse.nome}.
\end{definition}

\begin{unipdcode}[language=Java]
public class Coin {
    private String name;
    private double value;

    public Coin(String name, double value) {
        this.name = name;
        this.value = value;
    }

    public double getExchangeValue(double exchangeRate) {
        double converted = value * exchangeRate;
        return converted;
    }
}
\end{unipdcode}

Nei due assegnamenti del costruttore i parametri omonimi
\texttt{name} e \texttt{value} nascondono i campi: il prefisso
\texttt{this} distingue il membro della classe dal parametro.
Si consideri invece un costruttore che riceve un parametro
\texttt{aName}: l'istruzione \texttt{String name = aName;}
dichiarerebbe una variabile locale anziché inizializzare il campo
\texttt{name}, che resterebbe \texttt{null}. La dichiarazione
sarebbe valida e l'errore sarebbe logico. Un problema analogo
si verifica quando una variabile locale nasconde un campo statico.

Java non permette di dichiarare due variabili locali con lo
stesso nome in ambiti che si sovrappongono, mentre consente
lo stesso identificatore in blocchi \texttt{if} e \texttt{else}
distinti oppure in metodi differenti. Lo \emph{shadowing} di
un campo, invece, è lecito ma può rendere il codice ambiguo:
per questo i nomi e le qualificazioni devono essere scelti in
modo da rendere esplicito a quale dato si intende accedere.
